\documentclass[../main.tex]{subfiles}

\begin{document}

% -----------------------------
\section{Add operations inside the function}
% -----------------------------
Create a new file \path{step4_func_with_body.py} and import packages: 
\begin{lstlisting} [language=Python, caption={Adding imports}]
from xdsl.context import Context
from xdsl.ir import Region, Block
from xdsl.dialects.builtin import Builtin, ModuleOp, (*@\textbf{i32}@*)
from xdsl.dialects import func, arith
\end{lstlisting}

Build context once again and provide main function in the same way as in previous step.
\\Build module after that:
\begin{lstlisting} [language=Python, caption={Build module}]
def build_module_with_body(ctx: Context) -> ModuleOp:
    # Create constants: 1 and 2 as i32
    c1 = arith.ConstantOp.from_int_and_width(1, 32)
    c2 = arith.ConstantOp.from_int_and_width(2, 32)

    # Add them: %2 = arith.addi %0, %1 : i32
    add = arith.AddiOp(c1.result, c2.result)

    # Return the result
    ret = func.ReturnOp(add.result)

    # Build a block that contains all operations in order
    body_block = Block(ops=[c1, c2, add, ret])

    # Create a function @my_func() -> i32 with that body
    fn = func.FuncOp(
        "my_func",
        ([], [i32]),  # no arguments, one i32 result
    )
    fn.body = Region(body_block)

    # Put the function inside a module
    module = ModuleOp(ops=[fn])
    return module
\end{lstlisting}
Resulting IR should look like:
\begin{lstlisting} [language=bash, caption={Printed module}]
builtin.module {
  func.func @my_func() -> i32 {
    %0 = arith.constant 1 : i32
    %1 = arith.constant 2 : i32
    %2 = arith.addi %0, %1 : i32
    func.return %2 : i32
  }
}
\end{lstlisting}
At this point, we have:
\begin{itemize}
    \item a working context
    \item a working module
    \item a working func.func
    \item a working body block with real operations in xDSL
\end{itemize}

\end{document}
